%! Matrices as Functions — Unit 4, Lesson 4.13 — Group Activity
\documentclass[10pt]{article}
\usepackage{precalculus-article}
\usepackage{precalculus-boxes}
\newcommand{\TallMath}[1]{$\displaystyle #1\rule[-1.4em]{0pt}{3.2em}$}

\begin{document}

\pageheader{Unit 4, Lesson 4.13}{Group Activity}
\namepartnerperiod

\vspace{0.08in}

\begin{scenariobox}[Transformation Lab: Matrices as Functions]{sky}
\small
In this activity you will build, apply, and compose linear transformation matrices.
Work with your partner. Show all computations --- answers without work earn no credit.
\end{scenariobox}

\vspace{0.08in}

%----------------------------------------------------------------------
% Tier R
%----------------------------------------------------------------------
\begin{tcolorbox}[colback=white, colframe=black!40,
    title=\textbf{Tier R --- Reading the Transformation Matrix},
    fonttitle=\bfseries, arc=2mm, left=3mm, right=3mm, top=2mm, bottom=2mm]
\small

A linear transformation $L$ maps $\hat{\imath}=\langle1,0\rangle$ to $\langle2,0\rangle$
and maps $\hat{\jmath}=\langle0,1\rangle$ to $\langle0,2\rangle$.

\vspace{0.08in}
\textbf{(a)} Write the matrix $A$ associated with $L$.
(Remember: the columns are the images of $\hat{\imath}$ and $\hat{\jmath}$.)

\vspace{0.06in}
\hspace{1em}
$A = \begin{bmatrix}\blank{1.2cm}&\blank{1.2cm}\\\blank{1.2cm}&\blank{1.2cm}\end{bmatrix}$

\vspace{0.12in}
\textbf{(b)} Compute $L\bigl(\langle3,-1\rangle\bigr) = A\begin{bmatrix}3\\-1\end{bmatrix}$.
Show the multiplication.

\writelines{3}

\vspace{0.06in}
\textbf{(c)} Find $\det A$ and state what it tells you about the area scale factor of $L$.

\vspace{0.06in}
\hspace{1em} $\det A = $ \blank{2.5cm} \qquad Area scale factor: \blank{2.5cm}

\vspace{0.06in}
\hspace{1em} Interpretation: \blank{8cm}

\end{tcolorbox}

\vspace{0.1in}

%----------------------------------------------------------------------
% Tier A
%----------------------------------------------------------------------
\begin{tcolorbox}[colback=white, colframe=black!40,
    title=\textbf{Tier A --- The Rotation Matrix and Composition},
    fonttitle=\bfseries, arc=2mm, left=3mm, right=3mm, top=2mm, bottom=2mm]
\small

\textbf{(a)} Write the rotation matrix $R(\theta)$ for $\theta = 30°$
using exact values ($\cos30°=\dfrac{\sqrt{3}}{2}$, $\sin30°=\dfrac{1}{2}$).

\vspace{0.06in}
\hspace{1em}
$R(30°) = \begin{bmatrix}\blank{2.2cm}&\blank{2.2cm}\\[4pt]\blank{2.2cm}&\blank{2.2cm}\end{bmatrix}$

\vspace{0.12in}
\textbf{(b)} Write $R(60°)$ using exact values ($\cos60°=\dfrac{1}{2}$, $\sin60°=\dfrac{\sqrt{3}}{2}$).

\vspace{0.06in}
\hspace{1em}
$R(60°) = \begin{bmatrix}\blank{2.2cm}&\blank{2.2cm}\\[4pt]\blank{2.2cm}&\blank{2.2cm}\end{bmatrix}$

\vspace{0.12in}
\textbf{(c)} Compute the product $R(60°)\cdot R(30°)$ and show it equals $R(90°)$.
Use exact values throughout.

\writelines{5}

\vspace{0.06in}
$R(60°)\cdot R(30°) = \begin{bmatrix}\blank{1.2cm}&\blank{1.2cm}\\\blank{1.2cm}&\blank{1.2cm}\end{bmatrix}$
\qquad Compare with $R(90°) = \begin{bmatrix}0&-1\\1&0\end{bmatrix}$. Are they equal? \blank{1.5cm}

\end{tcolorbox}

\vspace{0.1in}

%----------------------------------------------------------------------
% Tier E
%----------------------------------------------------------------------
\begin{tcolorbox}[colback=white, colframe=black!40,
    title=\textbf{Tier E --- Composition Order and Inverse},
    fonttitle=\bfseries, arc=2mm, left=3mm, right=3mm, top=2mm, bottom=2mm]
\small

Let $L_1$ have matrix $A = \begin{bmatrix}2&0\\0&2\end{bmatrix}$
(scale by 2) and $L_2$ have matrix $B = \begin{bmatrix}0&-1\\1&0\end{bmatrix}$
($90°$ rotation).

\vspace{0.08in}
\textbf{(a)} Compute $L_2 \circ L_1$ (apply $L_1$ first, then $L_2$). Matrix: $BA$.

\writelines{3}
$BA = \begin{bmatrix}\blank{1.2cm}&\blank{1.2cm}\\\blank{1.2cm}&\blank{1.2cm}\end{bmatrix}$

\vspace{0.1in}
\textbf{(b)} Compute $L_1 \circ L_2$ (apply $L_2$ first, then $L_1$). Matrix: $AB$.

\writelines{3}
$AB = \begin{bmatrix}\blank{1.2cm}&\blank{1.2cm}\\\blank{1.2cm}&\blank{1.2cm}\end{bmatrix}$

\vspace{0.06in}
Are $BA$ and $AB$ equal? \blank{1.5cm} \quad What does this tell you about composition order? \blank{5cm}

\vspace{0.1in}
\textbf{(c)} Find $B^{-1}$ (the inverse of $L_2$) using the $2\times2$ inverse formula.
What rotation does $L_2^{-1}$ perform?

\writelines{3}
$B^{-1} = \begin{bmatrix}\blank{1.2cm}&\blank{1.2cm}\\\blank{1.2cm}&\blank{1.2cm}\end{bmatrix}$
\qquad $L_2^{-1}$ performs: \blank{5cm}

\end{tcolorbox}

\end{document}
